% Source: ACM_Templates.pdf, PDF pp. 17-18.
% Corrections and validation: docs/validation.md.
\subsection{Legendre 公式：阶乘中的质因子}
\label{algo:legendre}
\AlgorithmInfo
    {求素数 $p$ 在 $n!$ 的质因数分解中的指数。}
    {时间 $O(\log_p(n+1))$，额外空间 $O(1)$。}
    {$n\geq0$、$p$ 为素数；使用有符号 64 位非负整数。}

\subsubsection*{原理}
$1,2,\ldots,n$ 中每个 $p$ 的倍数至少贡献一个 $p$，每个 $p^2$ 的倍数再额外贡献一个，以此类推。因此
\[
    v_p(n!)=\sum_{j\geq1}\left\lfloor\frac n{p^j}\right\rfloor.
\]
代码直接反复令 $n\gets\lfloor n/p\rfloor$ 并累加，避免显式计算 $p^j$ 溢出。组合数中的指数为
\[
    v_p\binom nk=v_p(n!)-v_p(k!)-v_p((n-k)!).
\]
若要求最大 $t$ 使合数 $m^t\mid n!$，先分解 $m=\prod p_i^{\alpha_i}$，再取 $\min_i\lfloor v_{p_i}(n!)/\alpha_i\rfloor$；不能把合数直接当作 $p$ 代入。

\subsubsection*{使用教程}
\InlineCode{vp_fact(n,p)} 中 $n$ 为阶乘参数、$p$ 为素数，返回 $p$ 在 $n!$ 中的指数，类型为 \InlineCode{long long}；无需预处理，$n=0$ 时返回 0。
\begin{lstlisting}[style=mtycpp,numbers=none]
long long e=vp_fact(100,2);
\end{lstlisting}

\needspace{7\baselineskip}
\subsubsection*{模板代码}
\CodeFile{code/number-theory/legendre.cpp}
